% math.tex — verify-numerics experiment
% Problem specification, exact algorithm, and numerical error analysis.
%
% Part I — Problem, structure, exact algorithm:
%   def:hk2017-problem — full HK2017 capacity problem
%   fact:hk2017 — maximizer = minimum-action simple Reeb orbit
%   rem:min-support — minimum-support maximizers have beta > 0
%   rem:adjacency-pruning — directed transition prefilter
%   rem:qp-setup — per-sigma notation (A, F, V, beta_0, H', g, gamma_j)
%   prop:basic-properties — A may be empty, F is compact
%   prop:q-on-affine — Q restricted to affine plane
%   prop:critical-points — 3 cases for critical-point structure
%   prop:boundary-vs-interior — B1-B4 (boundary) and I1-I2 (interior)
%   alg:exact-solver — rational-arithmetic per-sigma solver
%   rem:exact-not-compute — what the solver skips (boundary projection)
%   rem:discontinuities — 6 discontinuities for f64
%
% Part II — Error analysis:
%   lem:q-error-first-order — first-order Q error bound via stationarity
%   lem:q-correction-second-order — Q correction reduces error to second order
%   rem:runtime-q-bound — computable runtime error bound
%   lem:pseudoinverse-orthogonality — structural property of pseudoinverse
%   cor:taylor-structure — first-order = -2x second-order
%   cor:exact-correction — correction exact in exact arithmetic
%
% ── Notational conventions ──
%
% Variables and their roles:
%   beta           free/generic variable ("maximize over beta")
%   beta_0         particular solution of C beta = d
%   beta^*         exact optimal beta (maximizer of Q on F)
%   \tilde\beta    f64-computed approximation of beta^*
%   \delta\beta    = \tilde\beta - beta^* (solution error)
%
% Decoration scheme (applies to all quantities, not just beta):
%   plain          exact mathematical value: H, C, V, H', g, gamma_j, Q
%   tilde (~)      f64-computed value: \tilde H, \tilde V, \tilde H', \tilde gamma_j
%   superscript *  optimal / solution: beta^*, Q^* = Q(beta^*)
%   subscript 0    particular / initial: beta_0
%   delta          perturbation / error: delta H = \tilde H - H
%
% Derived quantities (per-sigma, exact):
%   H, C, d        QP matrices (Definition def:hk2017-problem)
%   A              = {beta : C beta = d} (affine constraint set)
%   F              = A cap {beta >= 0} (closed feasible set)
%   V              orthonormal basis for ker(C), so A = beta_0 + im(V)
%   H'             = V^T H V (reduced Hessian on ker(C))
%   g              = V^T H beta_0 (gradient of Q|_A at alpha=0)
%   gamma_j        eigenvalues of H'
%   k              = dim ker(C) = m - rank(C)
%
% Standard reference: Higham, "Accuracy and Stability of Numerical
% Algorithms" (2002).  Plain = exact, tilde = computed follows Higham's
% convention.  Superscript * for optima is standard in optimization.

% \input'd from root math.tex — no \documentclass


% ══════════════════════════════════════════════════════════════════════
% Part I: Problem, structure, and exact algorithm
% ══════════════════════════════════════════════════════════════════════
%
% (1) HK2017 full problem
% (2) Simple Reeb orbit structure: minimum-support maximizers, pruning
% (3) General facts: quadratic forms on affine planes
% (4) Exact per-sigma solver
% (5) Continuity issues for f64


% ── 1. The HK2017 capacity problem ──────────────────────────────────

\begin{definition}[HK2017 capacity problem]\label{def:hk2017-problem}
Let $K = \{x \in \mathbb{R}^4 : a_i^\top x \leq 1,\; i = 1,\ldots,F\}$
be a convex polytope with dual vertices
$a_1, \ldots, a_F \in \mathbb{R}^4$.
For each subset $S \subseteq \{1, \ldots, F\}$ with $|S| = m \geq 2$
and each cyclic permutation $\sigma$ of~$S$, define:
\begin{itemize}
  \item \textbf{Action matrix}
    $H(\sigma) \in \mathbb{R}^{m \times m}$:
    for $i < j$,\;
    $H_{ij} = H_{ji} = \omega_0(a_{\sigma(i)},\, a_{\sigma(j)})$;\;
    $H_{ii} = 0$.
    (Symmetric by construction.
    $\omega_0$ is antisymmetric: $\omega_0(u, v) = -\omega_0(v, u)$.
    See Lemma~\ref{lem:H-quadratic}.)
  \item \textbf{Constraint matrix}
    $C(\sigma) \in \mathbb{R}^{5 \times m}$:
    rows $0$--$3$ are $C_{d,k} = a_{\sigma(k),d}$
    (closure: $\sum a_{\sigma(i)} \beta_i = 0$);
    row $4$ is all ones
    (normalization: $\sum \beta_i = 1$).
  \item $d = (0, 0, 0, 0, 1)^\top$.
\end{itemize}
The \emph{HK2017 capacity problem} is:
\begin{equation}\label{eq:hk2017-problem}
  Q_{\max}
  = \max_{\sigma}\;
    \max_{\beta}\;
    \tfrac{1}{2}\,\beta^\top H(\sigma)\, \beta
  \quad \text{s.t.} \quad
  C(\sigma)\beta = d, \;\; \beta \geq 0,
\end{equation}
where $\sigma$ ranges over all cyclic permutations of all subsets
$S \subseteq \{1, \ldots, F\}$ with $|S| \geq 2$.
The EHZ capacity is $c_{\mathrm{EHZ}}(K) = 1/(2\, Q_{\max})$.
\end{definition}

\begin{fact}[HK2017]\label{fact:hk2017}
% [TODO: JÖRN - cite HK2017 theorem precisely; which theorem number?]
Any global maximizer $(\sigma^*, \beta^*)$
of~\eqref{eq:hk2017-problem}
corresponds to a minimum-action simple closed Reeb orbit on
$\partial K$.
\end{fact}


% ── 2. Structure from Reeb orbits ────────────────────────────────────

\begin{remark}[Minimum-support maximizers]\label{rem:min-support}
A global maximizer $(\sigma^*, \beta^*)$
of~\eqref{eq:hk2017-problem}
may have $\beta^*_k = 0$ for some indices~$k$
(the orbit does not dwell on those facets).
Among all global maximizers achieving~$Q_{\max}$,
take one with minimum $|\operatorname{supp}(\beta^*)|$
(fewest nonzero components).
This minimum-support maximizer has $\beta^* > 0$:
if $\beta^*_k = 0$ for some~$k$, removing $\sigma(k)$
from the permutation gives a shorter $\sigma'$ with the same
$Q$ value (the $k$-th row/column of~$H$ contributes nothing,
and the constraints are satisfied by the remaining components).
This contradicts minimality of the support.

\textbf{Consequence:} the HK2017 algorithm can skip any
$(\sigma, \beta)$ where the maximum of~$Q$ over
$\{C(\sigma)\beta = d,\; \beta \geq 0\}$
is attained only on the boundary ($\beta_k = 0$ for some~$k$).
A shorter~$\sigma'$ achieves the same or higher~$Q$.
\end{remark}

\begin{remark}[Adjacency pruning]\label{rem:adjacency-pruning}
After the minimum-support reduction, one can further skip
permutations~$\sigma$ where consecutive pairs
$(\sigma(i), \sigma(i{+}1))$ fail the directed transition
feasibility test (Lemma~\ref{lem:numerical-transition-feasibility}):
no Reeb trajectory with that combinatorics exists.
This is independent of the $\beta$-optimization and is applied
as a prefilter before solving.
\end{remark}


% ── 3. Quadratic forms on affine planes ──────────────────────────────

% General facts, not specific to EHZ.

\begin{remark}[Setup and notation]\label{rem:qp-setup}
For a fixed~$\sigma$, write $H = H(\sigma)$,
$C = C(\sigma)$ (suppressing the $\sigma$-dependence).
Define:
\begin{itemize}
  \item $\mathcal{A} = \{\beta \in \mathbb{R}^m : C\beta = d\}$
    \quad (affine constraint set).
  \item $\mathcal{F} = \mathcal{A} \cap \{\beta \geq 0\}$
    \quad (closed feasible set).
  \item $\beta_0 \in \mathcal{A}$: a particular solution.
  \item $V \in \mathbb{R}^{m \times k}$: orthonormal basis for
    $\ker(C)$,\; $k = m - \operatorname{rank}(C)$,
    \;so $\mathcal{A} = \beta_0 + \operatorname{im}(V)$.
  \item $H' = V^\top H V \in \mathbb{R}^{k \times k}$:
    the \emph{reduced Hessian} (restriction of $H$ to~$\mathcal{A}$).
  \item $g = V^\top H \beta_0 \in \mathbb{R}^k$:
    the gradient of $Q|_{\mathcal{A}}$ at $\alpha = 0$.
  \item $\gamma_1, \ldots, \gamma_k$: eigenvalues of~$H'$.
\end{itemize}
The per-$\sigma$ problem is:
maximize $Q(\beta) = \tfrac{1}{2}\beta^\top H \beta$
over $\beta \in \mathcal{F}$.
\end{remark}

\begin{proposition}[Basic properties of $\mathcal{A}$ and $\mathcal{F}$]%
\label{prop:basic-properties}
\leavevmode
\begin{enumerate}
  \item $\mathcal{A}$ may be empty (if $C\beta = d$ is inconsistent).
  \item When $\mathcal{A} \neq \emptyset$,\;
    $\mathcal{F} = \mathcal{A} \cap \{\beta \geq 0\}$
    is compact: it is closed and bounded
    (contained in the simplex
    $\Delta = \{\beta \geq 0,\; \sum \beta_i = 1\}$
    by the normalization constraint).
    Hence $Q$ attains its maximum on~$\mathcal{F}$.
\end{enumerate}
\end{proposition}

\begin{proposition}[$Q$ on the affine plane]\label{prop:q-on-affine}
Parametrize $\mathcal{A}$ as $\beta = \beta_0 + V\alpha$,
$\alpha \in \mathbb{R}^k$.  Then
\begin{equation}\label{eq:q-restricted}
  Q(\beta_0 + V\alpha)
  = \underbrace{\tfrac{1}{2}\,\alpha^\top H' \alpha}_{\text{quadratic}}
  + \underbrace{g^\top \alpha}_{\text{linear}}
  + \underbrace{Q(\beta_0)}_{\text{constant}}.
\end{equation}
Setting $\nabla_\alpha = 0$ gives the stationarity condition
$H' \alpha + g = 0$.
\end{proposition}

\begin{proposition}[Critical-point structure]\label{prop:critical-points}
The critical points of $Q|_{\mathcal{A}}$
(solutions of $H'\alpha + g = 0$) form one of three cases:
\begin{enumerate}
  \item \textbf{$H'$ invertible} ($\det H' \neq 0$):
    unique critical point $\alpha^* = -(H')^{-1} g$,
    hence $\beta^* = \beta_0 + V\alpha^*$.
  \item \textbf{$H'$ singular, $g \perp \ker(H')$}
    (i.e., $g \in \operatorname{im}(H')$):
    affine subspace of critical points,
    $\alpha^*_{\mathrm{part}} + \ker(H')$,
    where $\alpha^*_{\mathrm{part}}$ is any particular solution.
    $Q$ takes the same value at all of them
    (Lemma~\ref{lem:well-defined}).
  \item \textbf{$H'$ singular, $g \not\perp \ker(H')$}:
    no critical points.
    $Q|_{\mathcal{A}}$ is affine (hence unbounded)
    along directions in $\ker(H')$ where $g$ has nonzero projection.
\end{enumerate}
\end{proposition}

\begin{proposition}[Boundary vs.\ interior maximum on $\mathcal{F}$]%
\label{prop:boundary-vs-interior}
Write $Q_{\mathcal{F}} = \max_{\beta \in \mathcal{F}} Q(\beta)$
and $\mathcal{F}_{\max} = \arg\max_{\beta \in \mathcal{F}} Q(\beta)
\subseteq \mathcal{F}$.
Define the boundary
$\partial\mathcal{F} = \{\beta \in \mathcal{F} : \beta_k = 0
\text{ for some } k\}$.

\medskip\noindent
\textbf{Cases where}
$\mathcal{F}_{\max} \subseteq \partial\mathcal{F}$
(max is on the boundary):
\begin{enumerate}
  \item[(B1)] Some $\gamma_j > 0$:
    $Q$ is unbounded on~$\mathcal{A}$
    (grows as $\tfrac{1}{2}\gamma_j t^2$ along eigenvector~$j$).
    Any interior maximum of $Q|_{\mathcal{A}}$ would require
    $H'$ negative semidefinite (second-order necessary condition),
    contradicting $\gamma_j > 0$.
    So $\mathcal{F}_{\max} \subseteq \partial\mathcal{F}$.

  \item[(B2)] All $\gamma_j \leq 0$, unique critical point
    (case~1 of Prop.~\ref{prop:critical-points}),
    but $\beta^*_k \leq 0$ for some~$k$:
    the global maximum of~$Q$ on~$\mathcal{A}$ is at~$\beta^*$
    (which is the unique critical point of a concave function),
    but $\beta^* \notin \{\beta > 0\}$,
    so $\mathcal{F}_{\max} \subseteq \partial\mathcal{F}$.

  \item[(B3)] $H'$ singular, $g \not\perp \ker(H')$
    (case~3 of Prop.~\ref{prop:critical-points}):
    $Q|_{\mathcal{A}}$ unbounded, same argument as~(B1).

  \item[(B4)] $H'$ singular, $g \perp \ker(H')$,
    critical subspace does not intersect~$\{\beta > 0\}$:
    all critical points have some $\beta_k \leq 0$,
    so $\mathcal{F}_{\max} \subseteq \partial\mathcal{F}$.
\end{enumerate}

\medskip\noindent
\textbf{Cases where}
$\mathcal{F}_{\max} \cap \{\beta > 0\} \neq \emptyset$
(max is at an interior point):
\begin{enumerate}
  \item[(I1)] All $\gamma_j < 0$ (negative definite $H'$),
    unique critical point $\beta^* > 0$:
    $\beta^*$ is the global maximum on~$\mathcal{A}$
    (concavity) and lies in the interior of~$\mathcal{F}$.
    $Q_{\mathcal{F}} = Q(\beta^*)$.

  \item[(I2)] $H'$ negative semidefinite, $g \perp \ker(H')$,
    critical subspace intersects $\{\beta > 0\}$:
    $Q_{\mathcal{F}} = Q(\beta^*)$ for any critical
    $\beta^* > 0$ in the subspace.
\end{enumerate}

In both~(I1) and~(I2), the interior maximizer has $\beta^* > 0$
and $Q(\beta^*)$ is the per-$\sigma$ contribution to~$Q_{\max}$.
\end{proposition}

\begin{proof}
(B1): If $\mathcal{F}_{\max} \not\subseteq \partial\mathcal{F}$,
some maximizer $\beta^* \in \mathcal{F}$ has $\beta^* > 0$.
Then $\beta^*$ is a local max of $Q$ on~$\mathcal{A}$
(locally $\mathcal{F}$ agrees with $\mathcal{A}$).
The second-order necessary condition requires $H'$ negative
semidefinite, contradicting $\gamma_j > 0$.
(The max may be interior to a \emph{face} of~$\mathcal{F}$
where some $\beta_k = 0$; this is still in~$\partial\mathcal{F}$.)

(B2): $H'$ negative definite means $Q|_{\mathcal{A}}$ is strictly
concave with unique maximum at~$\beta^*$.
For any $\beta \in \mathcal{F}$, $Q(\beta) \leq Q(\beta^*)$.
Since $\beta^* \notin \mathcal{F}^{\circ}$,
the max on~$\mathcal{F}$ is on~$\partial\mathcal{F}$.

(B3): Same as (B1) --- no critical point exists, so no interior max.

(B4): $Q$ is constant on the critical subspace.
Any $\beta \notin$ critical subspace has strictly lower~$Q$
(the non-null eigenvalues are $< 0$).
Since the critical subspace misses~$\{\beta > 0\}$,
the max on~$\mathcal{F}$ is on~$\partial\mathcal{F}$.

(I1), (I2): The critical $\beta^*$ achieves the global max of~$Q$
on~$\mathcal{A}$ (concavity / negative semidefiniteness).
Since $\beta^* > 0$, it is in~$\mathcal{F}^{\circ}$ and
$Q_{\mathcal{F}} = Q(\beta^*)$.
\end{proof}


% ── 4. Exact per-sigma solver ────────────────────────────────────────

\begin{algorithm}[Exact per-$\sigma$ solver]\label{alg:exact-solver}
Given exact $(H, C, d)$ over~$\mathbb{Q}$:

\medskip\noindent
\textbf{Step~1.}
Compute $\mathcal{A}$, $V$, $\beta_0$, $H'$, $g$, $\gamma_j$
(notation of Remark~\ref{rem:qp-setup}).
If $\mathcal{A} = \emptyset$: return DROP.

\medskip\noindent
\textbf{Step~2 (Boundary test).}
Determine whether $\mathcal{F}_{\max} \subseteq \partial\mathcal{F}$
using Proposition~\ref{prop:boundary-vs-interior}:
\begin{itemize}
  \item Check (B1): any $\gamma_j > 0$?
  \item Check (B3): $H'$ singular and $g \not\perp \ker(H')$?
  \item If $H'$ invertible:
    compute $\beta^* = \beta_0 - V(H')^{-1} g$.
    Check (B2): any $\beta^*_k \leq 0$?
  \item If $H'$ singular and $g \perp \ker(H')$:
    check (B4): does the critical subspace
    miss $\{\beta > 0\}$? (Linear feasibility in~$\gamma$.)
\end{itemize}

If any of (B1)--(B4) holds: return DROP.
By Remark~\ref{rem:min-support}, a shorter
permutation achieves the same~$Q$.
The HK2017 loop finds it.

\medskip\noindent
\textbf{Step~3 (Report).}
We are in case (I1) or (I2).
Return $Q_{\mathcal{F}} = Q(\beta^*)$
for any critical $\beta^*$ with $\beta^* > 0$.
\end{algorithm}

\begin{remark}[What the exact solver does \emph{not} compute]%
\label{rem:exact-not-compute}
When $\mathcal{F}_{\max} \subseteq \partial\mathcal{F}$ (cases B1--B4),
the solver returns DROP without computing $Q_{\mathcal{F}}$ or
$\mathcal{F}_{\max}$.
The complete answer would require projecting onto each boundary
face ($\beta_k = 0$) and optimizing there --- but the HK2017 loop
already does this by enumerating shorter permutations~$\sigma'$
(at least the ones that matter for~$Q_{\max}$).
\end{remark}


% ── 4b. Near-boundary dropping ───────────────────────────────────────

\begin{lemma}[Q loss from dropping a near-boundary component]%
\label{lem:near-boundary-drop}
Let $\beta^* \in \mathcal{F}$ achieve
$Q_{\mathcal{F}} = \max_{\beta \in \mathcal{F}} Q(\beta)$,
with $\beta^*_k \leq \varepsilon$ for some index~$k$
and $\beta^*_i \geq \delta > 0$ for all $i \neq k$.
Let $\sigma' = \sigma$ with $\sigma(k)$ removed,
with constraint matrix $C_{-k} \in \mathbb{R}^{5 \times (m-1)}$
(column~$k$ of~$C$ deleted), and feasible set
$\mathcal{F}' = \{C_{-k}\beta' = d,\; \beta' \geq 0\}$.

Assume $\sigma_{\min}(C_{-k}) > 0$ and
\begin{equation}\label{eq:positivity-condition}
  \delta > \frac{\|c_k\|}{\sigma_{\min}(C_{-k})} \cdot \varepsilon,
\end{equation}
where $c_k = C_{:,k}$ is the $k$-th column of~$C$.
Write $H_{-k} \in \mathbb{R}^{(m-1) \times (m-1)}$ for $H$
with row and column~$k$ deleted,
and $Q_{-k}(\beta') = \tfrac{1}{2}\beta'^\top H_{-k}\, \beta'$.

Then $Q^*(\sigma') \geq Q^*(\sigma) - E_{\mathrm{drop}}$ where
\begin{equation}\label{eq:combined-drop-bound}
  E_{\mathrm{drop}}
  = \varepsilon \left(
    \|H\| \cdot \|\beta^*\|
    + \frac{\|H\| \cdot \|\beta^*\| \cdot \|c_k\|}
           {\sigma_{\min}(C_{-k})}
    + \frac{\|H\| \cdot \|c_k\|^2 \varepsilon}
           {2\, \sigma_{\min}(C_{-k})^2}
  \right).
\end{equation}
\end{lemma}

\begin{proof}
Define $\beta' \in \mathbb{R}^{m-1}$ by deleting component~$k$
from~$\beta^*$ (reindexed to $\{1, \ldots, m-1\}$).

\emph{Step 1: Constraint violation of $\beta'$.}
$C_{-k}\beta' = \sum_{i \neq k} C_{:,i}\, \beta^*_i
= C\beta^* - c_k\, \beta^*_k = d - c_k\, \beta^*_k$.
So $\|C_{-k}\beta' - d\| = \|c_k\| \cdot \beta^*_k
\leq \|c_k\| \cdot \varepsilon$.

\emph{Step 2: Feasibility restoration.}
Define the projection onto $\{C_{-k}\beta' = d\}$:
$\bar\beta' = \beta' + C_{-k}^\top (C_{-k} C_{-k}^\top)^{-1}
  (d - C_{-k}\beta')$.
Then $C_{-k}\bar\beta' = d$ and
\[
  \|\bar\beta' - \beta'\|
  \leq \frac{\|c_k\|}{\sigma_{\min}(C_{-k})} \cdot \beta^*_k
  \leq \frac{\|c_k\|}{\sigma_{\min}(C_{-k})} \cdot \varepsilon.
\]
Componentwise: $\bar\beta'_i \geq \beta^*_i - \|\bar\beta' - \beta'\|
\geq \delta - \|c_k\| \varepsilon / \sigma_{\min}(C_{-k}) > 0$
by~\eqref{eq:positivity-condition}.
So $\bar\beta' \in \mathcal{F}'$ (feasible for~$\sigma'$).

\emph{Step 3: Q change from removing component~$k$.}
Expanding $Q(\beta^*) = \tfrac{1}{2}\sum_{i,j} H_{ij}\beta^*_i\beta^*_j$
and using $H_{kk} = 0$:
\begin{equation}\label{eq:q-removal}
  Q(\beta^*) - Q_{-k}(\beta')
  = \beta^*_k \sum_{j \neq k} H_{kj}\, \beta^*_j.
\end{equation}
Hence $|Q(\beta^*) - Q_{-k}(\beta')|
\leq \beta^*_k \cdot \|H_{k,:}\| \cdot \|\beta^*\|
\leq \varepsilon \cdot \|H\| \cdot \|\beta^*\|$
(Cauchy--Schwarz; $\|H_{k,:}\| \leq \|H\|$).

\emph{Step 4: Q change from the feasibility correction.}
By the triangle inequality on the quadratic Taylor expansion:
\[
  |Q_{-k}(\bar\beta') - Q_{-k}(\beta')|
  \leq \|H_{-k}\| \cdot \|\beta'\| \cdot \|\bar\beta' - \beta'\|
    + \tfrac{1}{2}\|H_{-k}\| \cdot \|\bar\beta' - \beta'\|^2.
\]
Since $\|H_{-k}\| \leq \|H\|$ (principal submatrix)
and $\|\beta'\| \leq \|\beta^*\|$:
\[
  |Q_{-k}(\bar\beta') - Q_{-k}(\beta')|
  \leq \frac{\|H\| \cdot \|\beta^*\| \cdot \|c_k\| \varepsilon}
            {\sigma_{\min}(C_{-k})}
    + \frac{\|H\| \cdot \|c_k\|^2 \varepsilon^2}
            {2\, \sigma_{\min}(C_{-k})^2}.
\]

\emph{Step 5: Combined bound.}
Since $Q^*(\sigma') \geq Q_{-k}(\bar\beta')$ (maximum $\geq$ any feasible value):
\begin{align*}
  Q^*(\sigma) - Q^*(\sigma')
  &\leq Q(\beta^*) - Q_{-k}(\bar\beta') \\
  &\leq |Q(\beta^*) - Q_{-k}(\beta')| + |Q_{-k}(\beta') - Q_{-k}(\bar\beta')|
  \leq E_{\mathrm{drop}}. \qedhere
\end{align*}
\end{proof}


\begin{remark}[Trinary treatment of $\beta > 0$]%
\label{rem:trinary-beta}
The test $\min_k \beta^*_k > 0$ is discontinuous.
The f64 solver computes $\tilde\beta$ and classifies:
\begin{itemize}
  \item $\min_k \tilde\beta_k > \varepsilon_\beta$
    $\;\Rightarrow\;$ \textsc{True}:
    by perturbation bounds,
    $\beta^*_k \geq \tilde\beta_k - \eta > 0$
    (where $\eta$ bounds $\|\tilde\beta - \beta^*\|_\infty$).
    KEEP: report $\tilde Q$.
  \item $\min_k \tilde\beta_k < -\varepsilon_\beta$
    $\;\Rightarrow\;$ \textsc{False}:
    $\beta^*_k < 0$, exact maximizer outside $\{\beta > 0\}$.
    DROP.
  \item $|\min_k \tilde\beta_k| \leq \varepsilon_\beta$
    $\;\Rightarrow\;$ \textsc{Indeterminate}:
    cannot determine sign of~$\beta^*_k$.
    \textbf{Also DROP}, justified by Lemma~\ref{lem:near-boundary-drop}:
\end{itemize}

\noindent
\textbf{Why Indeterminate $\to$ DROP is safe.}
If $|\min_k \tilde\beta_k| \leq \varepsilon_\beta$,
then $\beta^*_k \in [-\varepsilon_\beta - \eta,\;
\varepsilon_\beta + \eta]$.
Two sub-cases:
\begin{enumerate}
  \item $\beta^*_k \leq 0$: exact maximizer is on the boundary
    or outside. DROP is correct --- a shorter $\sigma'$ achieves
    $Q^*(\sigma') \geq Q^*(\sigma)$.
  \item $\beta^*_k > 0$ but $\beta^*_k \leq \varepsilon_\beta + \eta$:
    exact maximizer is interior but near-boundary.
    By Lemma~\ref{lem:near-boundary-drop} (with
    $\varepsilon = \varepsilon_\beta + \eta$),
    DROP loses at most $E_{\mathrm{drop}}$ in~$Q$.
\end{enumerate}
In both sub-cases, DROP is safe with bounded Q~loss.

\textbf{Critical assumption:} the capacity-achieving orbit
(minimum-support maximizer, Remark~\ref{rem:min-support})
has $\min_k \beta^*_k > \varepsilon_\beta + \eta$,
so it is classified \textsc{True} and never DROPped.
If this fails (margin below f64 resolution),
rational arithmetic fallback is needed.
\end{remark}


% ── 5. Continuity issues for f64 ────────────────────────────────────

\begin{remark}[Discontinuities in the exact algorithm]%
\label{rem:discontinuities}
Algorithm~\ref{alg:exact-solver} uses several tests that are
\textbf{discontinuous} functions of the input $(a_1, \ldots, a_F)$.
An f64 implementation cannot evaluate these tests exactly.
We list each discontinuity and what makes it problematic.

\begin{enumerate}
  \item \textbf{Sign of $\gamma_j$ (eigenvalues of $H'$).}
    Eigenvalue \emph{values} are continuous functions of~$H'$
    (Weyl's theorem), but $H'$ itself depends on $V$ (the
    null-space basis of~$C$), which changes discontinuously
    when the rank of~$C$ changes.
    Even at fixed rank, the sign test $\gamma_j > 0$ vs.\
    $\gamma_j \leq 0$ is discontinuous at $\gamma_j = 0$.
    Near $\gamma_j = 0$, f64 rounding can flip the sign.

  \item \textbf{Sign of $\beta^*_k$.}
    $\beta^* = \beta_0 - V(H')^{-1} g$ involves $(H')^{-1}$,
    which amplifies errors by $1/|\gamma_j|$ in eigendirection~$j$.
    When $|\gamma_j|$ is small, $\beta^*$ is ill-conditioned:
    small perturbations of $a_i$ cause large changes in~$\beta^*$.
    The test $\beta^*_k > 0$ vs.\ $\beta^*_k \leq 0$ is
    discontinuous at $\beta^*_k = 0$, and near this boundary,
    the ill-conditioning of~$\beta^*$ makes the test unreliable.

  \item \textbf{Rank of $C$.}
    The dimension $k = m - \operatorname{rank}(C)$
    of the null space is an integer-valued function of~$C$.
    A small perturbation can change rank, which changes the
    dimension of~$\mathcal{A}$ and the size of~$H'$.
    In f64, rank is determined by thresholding singular values
    of~$C$, which is discontinuous at the threshold.

  \item \textbf{Rank of $H'$ (semidefinite case).}
    The solvability check $g \perp \ker(H')$
    (Prop.~\ref{prop:critical-points}, case~2 vs.\ case~3)
    depends on the rank of~$H'$.
    Same discontinuity as item~3.

  \item \textbf{Feasibility of the critical subspace (case B4).}
    The LP ``does the critical subspace intersect
    $\{\beta > 0\}$?'' has a binary answer that is discontinuous
    in the LP coefficients, which depend on~$V$, $H'$, $g$.

  \item \textbf{Eigenvalue threshold for the reduced Hessian.}
    An f64 solver must choose a threshold $\varepsilon$ separating
    ``clearly negative'' from ``near-zero'' eigenvalues.
    Eigenvectors with $|\gamma_j| > \varepsilon$ contribute to
    $(H')^{-1}$ (and hence to $\beta^*$); those with
    $|\gamma_j| \leq \varepsilon$ are treated as null
    (and contribute to the search space for $\beta > 0$).
    This inclusion/exclusion is a discontinuous function of
    $\gamma_j$ at $\gamma_j = \pm\varepsilon$.
    At the threshold, the eigenvector's contribution to
    $\beta^*$ is $O(g_j / \varepsilon)$, which can be large.

    % [GAP - Need an argument that at the threshold, the
    % contribution to Q is small (it should be, since Q varies
    % as O(gamma_j * |delta_beta_j|^2) near critical points,
    % and gamma_j = epsilon is small).  But the contribution
    % to the beta > 0 test may NOT be small.]
\end{enumerate}
\end{remark}


% ══════════════════════════════════════════════════════════════════════
% Part II: Error analysis
% ══════════════════════════════════════════════════════════════════════
%
% Structure:
%   (A) Perturbation chain: a_i -> C -> V -> H' -> gamma_j -> beta
%   (B) Q error bounds (existing, from earlier work)
%   (C) beta > 0 certification (the main open problem)


% ── A. Perturbation chain ────────────────────────────────────────────

\begin{remark}[Perturbation chain overview]\label{rem:perturbation-chain}
The f64 solver computes a chain of quantities, each depending
on the previous.
At each link, errors from f64 arithmetic and from perturbed
inputs accumulate.
We bound the error at each link using standard backward
stability results for the f64 primitives (SVD, eigendecomposition)
and perturbation theory (Weyl, Davis--Kahan).

The chain, with the key bound at each link:
\[
  a_i
  \;\xrightarrow{\text{assemble}}\;
  C, H
  \;\xrightarrow{\text{SVD}}\;
  V, \beta_0
  \;\xrightarrow{V^\top H V}\;
  H'
  \;\xrightarrow{\text{eig}}\;
  \gamma_j
  \;\xrightarrow{(H')^{-1} g}\;
  \alpha^*
  \;\xrightarrow{\beta_0 + V\alpha^*}\;
  \beta^*.
\]
\end{remark}


\begin{lemma}[Link 1: Assembly]\label{lem:link-assembly}
$H$ and $C$ are assembled from the dual vertices $a_i$ by
$O(m^2)$ elementary f64 operations
($\omega_0$ evaluations, copies).
The computed $\tilde C$ and $\tilde H$ satisfy
\[
  \|\tilde C - C\| \leq c_1\, \varepsilon_{\mathrm{mach}}\, \|C\|,
  \qquad
  \|\tilde H - H\| \leq c_2\, \varepsilon_{\mathrm{mach}}\, \|H\|,
\]
where $c_1, c_2$ are small constants depending on~$m$
(at most $O(m)$), and $\varepsilon_{\mathrm{mach}} \approx 1.1 \times 10^{-16}$.
% [TODO: JÖRN - verify constants; the assembly involves
% a subtraction (omega_0 = q1*p2 - q2*p1) which could cause
% cancellation if a_i are nearly aligned.  For our polytopes,
% a_i are O(1) with no near-cancellation, so this is safe.]

In our setting, $\|C\|$ and $\|H\|$ are $O(1)$
(dual vertices $a_i$ are $O(1)$),
so $\|\delta C\|$ and $\|\delta H\|$ are $O(\varepsilon_{\mathrm{mach}})$.
\end{lemma}


\begin{lemma}[Link 2: Null-space basis via SVD]\label{lem:link-svd}
% [TODO: JÖRN - cite specific theorem from Higham or GVL
% for backward stability of SVD and null-space computation.
% Standard result: computed SVD is exact SVD of C + delta C
% with ||delta C|| <= c * eps_mach * ||C||.]
The SVD of~$\tilde C$ is backward stable: the computed
$\tilde U \tilde\Sigma \tilde V_{\mathrm{full}}^\top$
is the exact SVD of $\tilde C + \delta C_{\mathrm{SVD}}$
where $\|\delta C_{\mathrm{SVD}}\| \leq c_3\, \varepsilon_{\mathrm{mach}}\, \|\tilde C\|$.

The null-space basis $\tilde V$ (the last $k$ right singular vectors)
approximates the exact null space of~$C$.
The angle between the computed and exact null spaces is bounded by
the \textbf{Davis--Kahan $\sin\theta$ theorem}:
\[
  \|\sin\Theta(\tilde V, V)\|
  \leq \frac{\|\delta C_{\mathrm{total}}\|}
       {\sigma_p(C) - \sigma_{p+1}(C)}
  = \frac{\|\delta C_{\mathrm{total}}\|}
       {\sigma_{\min}(C)},
\]
where $\delta C_{\mathrm{total}}$ combines assembly and SVD errors,
$\sigma_p(C) = \sigma_{\min}(C)$ is the smallest nonzero singular
value ($C$ has full row rank~$p = 5$),
and $\sigma_{p+1}(C) = 0$ (the gap to the null space).

% [GAP - The sin(Theta) bound requires sigma_min(C) >> 0.
% If sigma_min(C) is small, the null space of C is
% ill-conditioned: a small perturbation of C changes V
% significantly.  In the EHZ setting, sigma_min(C) depends
% on the geometry of the dual vertices a_{sigma(i)}.
% Empirically, sigma_min(C) < 0.01 is associated with
% boundary-case beta (margin ~ 1e-12), so ill-conditioned
% null spaces coincide with cases the solver should return
% INDETERMINATE anyway.]
\end{lemma}


\begin{lemma}[Link 3: Reduced Hessian]\label{lem:link-reduced-hessian}
The reduced Hessian $H' = V^\top H V$ is computed as
$\tilde H' = \tilde V^\top \tilde H\, \tilde V$.
Write $\tilde V = V + \delta V$ where $\|\delta V\|$ is bounded
by Lemma~\ref{lem:link-svd}. Then:
\begin{align*}
  \tilde H'
  &= (V + \delta V)^\top (H + \delta H) (V + \delta V) \\
  &= V^\top H V
     + V^\top H\, \delta V
     + \delta V^\top H V
     + V^\top \delta H\, V
     + O(\|\delta V\|^2, \|\delta H\| \cdot \|\delta V\|) \\
  &= H' + \Delta H',
\end{align*}
where the first-order perturbation satisfies
\[
  \|\Delta H'\|
  \leq 2\,\|H\|\,\|\delta V\| + \|\delta H\| + O(\text{second order}).
\]
Using $\|\delta H\| = O(\varepsilon_{\mathrm{mach}})$
and $\|\delta V\| = O(\varepsilon_{\mathrm{mach}} / \sigma_{\min}(C))$
from Lemmas~\ref{lem:link-assembly}--\ref{lem:link-svd}:
\[
  \|\Delta H'\|
  = O\!\left(
    \frac{\|H\|\, \varepsilon_{\mathrm{mach}}}{\sigma_{\min}(C)}
  \right).
\]
\end{lemma}


\begin{lemma}[Link 4: Eigenvalues of the reduced Hessian]%
\label{lem:link-eigenvalues}
% [TODO: JÖRN - cite Weyl's theorem precisely]
By Weyl's eigenvalue perturbation theorem,
the computed eigenvalues $\tilde\gamma_j$ of~$\tilde H'$
satisfy
\[
  |\tilde\gamma_j - \gamma_j|
  \leq \|\tilde H' - H'\|
  \leq \|\Delta H'\| + c_4\, \varepsilon_{\mathrm{mach}}\, \|\tilde H'\|,
\]
where the second term accounts for f64 rounding in the
eigendecomposition itself (backward stable).
Combining with Lemma~\ref{lem:link-reduced-hessian}:
\begin{equation}\label{eq:eigenvalue-perturbation}
  |\tilde\gamma_j - \gamma_j|
  = O\!\left(
    \frac{\|H\|\, \varepsilon_{\mathrm{mach}}}{\sigma_{\min}(C)}
  \right).
\end{equation}

\textbf{Eigenvalue sign reliability:}
if $|\tilde\gamma_j| > \varepsilon_\gamma$
where
$\varepsilon_\gamma = c_5 \cdot \|H\|\, \varepsilon_{\mathrm{mach}} / \sigma_{\min}(C)$,
then $\operatorname{sign}(\tilde\gamma_j) = \operatorname{sign}(\gamma_j)$.
The trinary classification threshold is derived, not heuristic.
\end{lemma}


\begin{lemma}[Link 5: Critical point $\beta^*$]\label{lem:link-beta}
In the negative definite case (all $\gamma_j < 0$),
the exact critical point is
$\beta^* = \beta_0 + V \alpha^*$ where
$\alpha^* = -(H')^{-1} g$.
The computed $\tilde\alpha = -(\tilde H')^{-1} \tilde g$
satisfies
\[
  \|\tilde\alpha - \alpha^*\|
  \leq \|(H')^{-1}\| \cdot
    \bigl(
      \|\Delta H'\| \cdot \|\alpha^*\| + \|\delta g\|
    \bigr)
    + O(\text{second order}),
\]
% [TODO: JÖRN - this is the standard perturbation bound for
% linear solves: (A+dA)^{-1}(b+db) - A^{-1}b.  Cite.]
where $\|(H')^{-1}\| = 1/\min_j |\gamma_j|$.

\textbf{Direction-dependent amplification:}
In the eigenbasis of $H'$, the $j$-th component of
$\tilde\alpha - \alpha^*$ is amplified by $1/|\gamma_j|$.
For small $|\gamma_j|$, the error in direction~$j$ is large.
Consequently:
\[
  \tilde\beta_k - \beta^*_k
  = \sum_j (V w_j)_k \cdot
    \frac{\delta(\text{stuff})_j}{|\gamma_j|},
\]
where $w_j$ are the eigenvectors of~$H'$.
When $|\gamma_j|$ is small, component~$k$ of $\tilde\beta$
can differ from $\beta^*_k$ by much more than
$\varepsilon_{\mathrm{mach}}$.

\textbf{$\beta > 0$ certification:}
We need $\min_k \beta^*_k > 0$.  The margin test is:
\[
  \min_k \tilde\beta_k
  > \sum_j \frac{|(\text{error bound})_j|}{|\gamma_j|}
    \cdot \max_k |(V w_j)_k|.
\]
If this holds, $\beta^* > 0$ is certified.
If not, return INDETERMINATE.

% [GAP - The bound above is schematic, not rigorous.
% The "delta(stuff)_j" involves both Delta H' and delta g
% projected onto eigenvector j.  Need to make this precise.
% Also: the bound requires knowing |gamma_j| (exact), but
% we only have |tilde gamma_j|.  By Lemma lem:link-eigenvalues,
% |gamma_j| >= |tilde gamma_j| - epsilon_gamma, so we can
% substitute -- but this fails when |tilde gamma_j| <= epsilon_gamma
% (the INDETERMINATE zone for eigenvalues).]
\end{lemma}


% ── B. Q error bounds (existing) ─────────────────────────────────────

% The lemmas below were proven before the perturbation chain above.
% They bound |Q(tilde beta) - Q(beta^*)| using the KKT residual r,
% without going through the chain.  This is a complementary approach:
% the chain gives componentwise understanding, while the residual
% bound gives a single scalar bound on Q error.
%
% The KKT system M x = b where M = [H, C^T; C, 0] is the augmented
% form of the per-sigma problem.  The residual r = M tilde x - b
% measures how well the computed solution satisfies the KKT conditions.

\begin{lemma}[First-order Q error bound]\label{lem:q-error-first-order}
Let $Q(\beta) = \tfrac{1}{2}\beta^\top H \beta$ where
$H \in \mathbb{R}^{m \times m}$ is symmetric.
Let $C \in \mathbb{R}^{p \times m}$ have full row rank,
$d \in \mathbb{R}^p$, and
$M = \bigl[\begin{smallmatrix} H & C^\top \\ C & 0 \end{smallmatrix}\bigr]$.
Suppose $\beta^* > 0$ solves the interior KKT system $M(\beta^*,\lambda^*)^\top = (0,d)^\top$,
and $(\tilde\beta, \tilde\lambda)$ is a computed solution with residual
$r = M(\tilde\beta,\tilde\lambda)^\top - (0,d)^\top$.
Write $\delta\beta = \tilde\beta - \beta^*$ and
$r_\lambda = C\tilde\beta - d = C\,\delta\beta$ for the
constraint residual (the last $p$ components of~$r$).
Then
\[
  |Q(\tilde\beta) - Q(\beta^*)|
  \;\leq\;
  \frac{\|H\|\,\|\beta^*\|}{\sigma_{\min}(C)}\,\|r\|
  \;+\; \tfrac{1}{2}\,\|H\|\,\|\delta\beta\|^2,
\]
where $\|\cdot\|$ denotes the spectral norm for matrices and the
Euclidean norm for vectors, and $\sigma_{\min}(C)$ is the smallest
singular value of~$C$.
\end{lemma}

\begin{proof}
Taylor expansion at $\beta^*$:
\begin{equation}\label{eq:q-taylor}
  Q(\tilde\beta) - Q(\beta^*)
  = (H\beta^*)^\top \delta\beta + \tfrac{1}{2}\,\delta\beta^\top H\,\delta\beta.
\end{equation}

\emph{First-order term.}
Since $\beta^* > 0$, the interior KKT stationarity condition gives
$H\beta^* + C^\top\lambda^* = 0$, i.e.,
$H\beta^* = -C^\top\lambda^*$. Substituting:
\begin{equation}\label{eq:first-order-rewrite}
  (H\beta^*)^\top\delta\beta
  = -(C^\top\lambda^*)^\top\delta\beta
  = -\lambda^{*\top} C\,\delta\beta
  = -\lambda^{*\top} r_\lambda.
\end{equation}
Since $C$ has full row rank~$p$, $C^\top \in \mathbb{R}^{m \times p}$
is injective (its null space is trivial). Therefore
$\|C^\top y\| \geq \sigma_{\min}(C^\top)\|y\| = \sigma_{\min}(C)\|y\|$
for all $y \in \mathbb{R}^p$. Applying this to $C^\top\lambda^* = -H\beta^*$:
\[
  \sigma_{\min}(C)\,\|\lambda^*\|
  \leq \|C^\top\lambda^*\|
  = \|H\beta^*\|
  \leq \|H\|\,\|\beta^*\|.
\]
Thus $\|\lambda^*\| \leq \|H\|\,\|\beta^*\| / \sigma_{\min}(C)$.
Since $r_\lambda$ is a sub-vector of $r$ (Euclidean norm),
$\|r_\lambda\| \leq \|r\|$.
Therefore
\[
  |(H\beta^*)^\top\delta\beta|
  = |\lambda^{*\top} r_\lambda|
  \leq \|\lambda^*\|\,\|r_\lambda\|
  \leq \frac{\|H\|\,\|\beta^*\|}{\sigma_{\min}(C)}\,\|r\|.
\]

\emph{Second-order term.}
$|\tfrac{1}{2}\,\delta\beta^\top H\,\delta\beta|
  \leq \tfrac{1}{2}\,\|H\|\,\|\delta\beta\|^2$.

Combining with \eqref{eq:q-taylor} by the triangle inequality:
\[
  |Q(\tilde\beta) - Q(\beta^*)|
  \leq \frac{\|H\|\,\|\beta^*\|}{\sigma_{\min}(C)}\,\|r\|
  + \tfrac{1}{2}\,\|H\|\,\|\delta\beta\|^2. \qedhere
\]
\end{proof}


\begin{lemma}[Q correction is second-order]\label{lem:q-correction-second-order}
Under the same assumptions as Lemma~\ref{lem:q-error-first-order},
define the \emph{Q correction} as $\Delta Q = \tilde\lambda^\top r_\lambda$,
where $\tilde\lambda$ is the computed Lagrange multiplier and
$r_\lambda = C\tilde\beta - d$.
Then the corrected Q value
$Q_{\mathrm{corr}} = Q(\tilde\beta) + \Delta Q$ satisfies
\[
  Q_{\mathrm{corr}} - Q(\beta^*)
  = \delta\lambda^\top r_\lambda
    + \tfrac{1}{2}\,\delta\beta^\top H\,\delta\beta,
\]
where $\delta\lambda = \tilde\lambda - \lambda^*$.
In particular,
$|Q_{\mathrm{corr}} - Q(\beta^*)|
  \leq \|\delta\lambda\|\,\|r_\lambda\|
    + \tfrac{1}{2}\,\|H\|\,\|\delta\beta\|^2$.
\end{lemma}

\begin{proof}
From \eqref{eq:q-taylor} and \eqref{eq:first-order-rewrite}:
$Q(\tilde\beta) - Q(\beta^*) = -\lambda^{*\top} r_\lambda
+ \tfrac{1}{2}\,\delta\beta^\top H\,\delta\beta$. Adding the correction:
\begin{align*}
  Q_{\mathrm{corr}} - Q(\beta^*)
  &= Q(\tilde\beta) + \tilde\lambda^\top r_\lambda - Q(\beta^*) \\
  &= -\lambda^{*\top} r_\lambda + \tilde\lambda^\top r_\lambda
    + \tfrac{1}{2}\,\delta\beta^\top H\,\delta\beta \\
  &= (\tilde\lambda - \lambda^*)^\top r_\lambda
    + \tfrac{1}{2}\,\delta\beta^\top H\,\delta\beta \\
  &= \delta\lambda^\top r_\lambda
    + \tfrac{1}{2}\,\delta\beta^\top H\,\delta\beta.
\end{align*}
The bound follows from the triangle inequality and
$|\delta\beta^\top H\,\delta\beta| \leq \|H\|\,\|\delta\beta\|^2$.
\end{proof}


\begin{remark}[Runtime error bound]\label{rem:runtime-q-bound}
At runtime, the exact $\beta^*$ is unknown.
The first-order term of Lemma~\ref{lem:q-error-first-order}
can be computed using the solver output $\tilde\beta$
as a proxy for~$\beta^*$:
\[
  E_1 = \frac{\|H\|\,\|\tilde\beta\|}{\sigma_{\min}(C)}\,\|r\|.
\]
This bounds the first-order part of the Q~error.
The second-order term $\tfrac{1}{2}\|H\|\,\|\delta\beta\|^2$
cannot be computed directly (since $\delta\beta$ is unknown),
but is negligible when $\|\delta\beta\| \ll \|\beta^*\|$,
which holds when the solver achieves small residuals on
well-conditioned problems.

% [GAP - The claim "||delta_beta|| is small when ||r|| is small" requires
% bounding ||delta_beta|| in terms of ||r|| via the conditioning of M.
% The naive bound ||delta_beta|| <= ||M^+|| * ||r|| is valid but
% ||M^+|| can be very large (up to 10^16 in experiments). A tighter
% analysis via the block structure of M is needed to make this rigorous.]

All quantities in $E_1$ are available from the solver:
$\|H\|$ from the eigendecomposition (largest $|\lambda_i(H)|$),
$\|\tilde\beta\|$ from the solution vector,
$\|r\|$ from the KKT residual, and
$\sigma_{\min}(C)$ from the SVD of~$C$
(already computed by the constraint solver).

% [TODO: JÖRN - Empirical validation shows max ratio
%   |Q_err| / E_1 = 0.196 across 477 test problems from 15 matrix families.
%   The constant 1 in the bound is not tight; the proof gives an upper
%   bound, not an equality.  Whether a tighter constant can be proven
%   (e.g., involving dimension m) is open.]
\end{remark}


\begin{lemma}[Pseudoinverse orthogonality]\label{lem:pseudoinverse-orthogonality}
Let $M$ be symmetric and $\tilde x = M^+ b$ the pseudoinverse solution.
Let $x^*$ satisfy $M x^* = b$ with $x^* \in \mathrm{col}(M)$.
Define $\delta x = \tilde x - x^*$ and $r = M\tilde x - b$.
Then $\delta x^\top M\, \delta x = 0$.
\end{lemma}

\begin{proof}
Since $M$ is symmetric, $M^+$ maps into $\mathrm{col}(M) = \mathrm{row}(M)$,
so $\tilde x \in \mathrm{col}(M)$. By assumption $x^* \in \mathrm{col}(M)$,
hence $\delta x \in \mathrm{col}(M)$.
From $M\tilde x = b + r$ and $M x^* = b$, we get $M\,\delta x = r$.
The residual satisfies $r \perp \mathrm{col}(M)$ (standard pseudoinverse property).
Therefore $\delta x^\top r = 0$ (since $\delta x \in \mathrm{col}(M)$), and
\[
  \delta x^\top M\,\delta x = \delta x^\top r = 0. \qedhere
\]
\end{proof}


\begin{corollary}[Taylor structure of Q error]\label{cor:taylor-structure}
Under the assumptions of Lemma~\ref{lem:q-error-first-order} and
Lemma~\ref{lem:pseudoinverse-orthogonality}, the uncorrected Q~error satisfies
\[
  Q(\tilde\beta) - Q(\beta^*)
  = (H\beta^*)^\top\delta\beta + \tfrac{1}{2}\,\delta\beta^\top H\,\delta\beta
\]
with $(H\beta^*)^\top\delta\beta = -\delta\beta^\top H\,\delta\beta$,
i.e., the first-order term is exactly $-2$ times the second-order term.
Consequently, $Q(\tilde\beta) - Q(\beta^*) = -\tfrac{1}{2}\,\delta\beta^\top H\,\delta\beta$.
\end{corollary}

\begin{proof}
Expanding $\delta x^\top M\,\delta x = 0$ in blocks:
\[
  \delta\beta^\top H\,\delta\beta + 2\,\delta\beta^\top C^\top\delta\lambda = 0.
\]
From $x^{*\top} M\,\delta x = 0$ (same argument with $x^*$ in place of $\delta x$):
\[
  \beta^{*\top} H\,\delta\beta + d^\top\delta\lambda + \lambda^{*\top} r_\lambda = 0,
\]
using $C\beta^* = d$ and $C\,\delta\beta = r_\lambda$.
Subtracting gives
$(H\beta^*)^\top\delta\beta - \delta\beta^\top H\,\delta\beta
= d^\top\delta\lambda + \lambda^{*\top} r_\lambda - 2\,r_\lambda^\top\delta\lambda$.
% [GAP - The identity (Hβ*)^T δβ = -δβ^T H δβ is verified empirically
% (ratio = 2.0000 across 95 rank-deficient cases) but the algebraic
% proof above yields extra terms d^T δλ + λ*^T r_λ - 2 r_λ^T δλ
% that need to cancel. This cancellation follows from x*^T r = 0
% (since x* ∈ col(M)), giving β*^T r_β + λ*^T r_λ = 0, combined
% with the block expansion. Full derivation needs Jörn.]
\end{proof}


\begin{corollary}[Exact correction in exact arithmetic]\label{cor:exact-correction}
Under the same assumptions,
\[
  Q_{\mathrm{corr}} - Q(\beta^*)
  = \delta\lambda^\top r_\lambda + \tfrac{1}{2}\,\delta\beta^\top H\,\delta\beta.
\]
From Lemma~\ref{lem:pseudoinverse-orthogonality},
$\tfrac{1}{2}\,\delta\beta^\top H\,\delta\beta = -r_\lambda^\top\delta\lambda$
(expanding $\delta x^\top M\,\delta x = 0$).
Since $\delta\lambda^\top r_\lambda$ and $r_\lambda^\top\delta\lambda$ are
both scalars and equal, $Q_{\mathrm{corr}} - Q(\beta^*) = 0$.
\end{corollary}
